\section{Equal-budget sequential all-pairs comparison}
\label{app:ustat_extension}
This extension uses the same compliance--success--cost generator and
three component thresholds as the primary study, but a separate seed
suite: seed 20260918 rather than the primary 20260917. It includes eight
scenarios with 2,000 repetitions each, plus four null/boundary scenarios
with 10,000 repetitions each. The main comparison uses 199 looks from
100 to 10,000 records per arm; group-sequential rules use ten planned
looks at equally spaced information fractions. A look with $n$ records
per arm costs $2n$ executions for either the disjoint or all-pairs method.
The saved disjoint-betting rows reproduce the separate online-comparator
pipeline exactly, not the primary table's different random seed stream.

\paragraph{Procedures and validity classes.}
The all-pairs estimate is $U_n=n^{-2}\sum_{i,j}h(A_i,B_j)$, with
row/column projection variance estimate
$(\widehat\sigma_A^2+\widehat\sigma_B^2)/n$.
Fixed Wald and the O'Brien--Fleming (OBF), Pocock, and Hwang--Shih--DeCani
(HSD, parameter $-3$) spending comparisons use the Gaussian-limit
covariance at nested looks, following the sequential win-statistic
frameworks of \citet{zhang2024sequential,bergemann2026group}.
The information fraction $n/N$ is a first-order approximation; the
nonlinear statistic does not have exactly independent Gaussian increments
at finite sample sizes. At the first look $n=100$, saved Monte Carlo
estimates of population variance components put the omitted residual term between
$5.94\times10^{-6}$ and $2.49\times10^{-5}$, or 0.13--0.41\% of the
first-order term. This is not a bound on plug-in variance-estimation error.

The projection-Gaussian reference subtracts a one-sided normal-mixture
boundary multiplied by the estimated projection standard deviation.
Its justification is asymptotic, using the time-uniform Gaussian
approximation perspective of \citet{waudbysmith2024timeuniform}.
For $X_i=(A_i,B_i)$ and
$k(X_i,X_j)=\{h(A_i,B_j)+h(A_j,B_i)\}/2$, its one-sample U-statistic
$U_n^*$ obeys $|U_n-U_n^*|\leq2/n$. Its linear projection term has
asymptotic variance $\sigma_A^2+\sigma_B^2$ under independent arms;
the first Hoeffding projection of $k$ itself has one quarter of that
variance. Bounded kernels,
nondegenerate projection variance and a consistent variance estimate
support the asymptotic comparison. This is not a finite-sample
5\% crossing guarantee from $n=100$, nor an implementation of the
delayed-start family in \citet{cai2026ustatistics}.
Disjoint betting retains its stated finite-sample guarantee.

\paragraph{Compare the same decision event.}
The primary comparator requires every gate to pass at the same look.
A separately reported retained-crossing rule remembers each gate's
first crossing. These are distinct procedures even under stationarity;
retention is not used to certify a current drifting conjunction.
Table~\ref{tab:ustat_calibration} separates component rejection from
the two guarded events at the compliance boundary $\Delta_C=-0.01$,
with the other targets favorable. All intervals are pointwise 95\%
Wilson Monte Carlo intervals. A rate compatible with 5\% at this
precision is not a proof of uniform calibration.

\begin{table}[h]\centering\scriptsize
\caption{Rare-compliance boundary: rates (\%) and Monte Carlo intervals,
10,000 repetitions. The three columns concern different events.}
\label{tab:ustat_calibration}
\begin{tabular}{lrrr}\toprule
Rule & Component gate & Guarded, same look & Guarded, retained\\\midrule
All-pairs fixed Wald & 5.51 [5.08, 5.97] & 5.50 [5.07, 5.96] & 5.50 [5.07, 5.96] \\
All-pairs OBF & 5.46 [5.03, 5.92] & 5.40 [4.97, 5.86] & 5.45 [5.02, 5.91] \\
All-pairs Pocock & 6.43 [5.97, 6.93] & 5.09 [4.68, 5.54] & 6.43 [5.97, 6.93] \\
All-pairs HSD & 5.67 [5.23, 6.14] & 5.23 [4.81, 5.68] & 5.66 [5.22, 6.13] \\
Projection Gaussian & 7.25 [6.76, 7.77] & 1.83 [1.59, 2.11] & 7.11 [6.62, 7.63] \\
Disjoint OBF & 5.45 [5.02, 5.91] & 5.39 [4.96, 5.85] & 5.44 [5.01, 5.90] \\
Disjoint betting & 0.76 [0.61, 0.95] & 0.31 [0.22, 0.44] & 0.72 [0.57, 0.91] \\
\bottomrule\end{tabular}\end{table}

The projection-Gaussian component excess does not contradict its
asymptotic validity statement, and the low guarded rate does not certify
the underlying component test. An independently reproduced diagnostic
finds that 410 of its 725 component rejections first occur by 500 records
per arm, with 239 having no observed noncompliant A records at crossing.
These conditional-on-rejection summaries are consistent with unstable
rare-event variance estimation, but do not prove the mechanism or a remedy.

\begin{table}[h]\centering\small
\caption{Same-look guarded decisions, separate 2,000-repetition comparison.
Use is mean capped records per arm, with its Monte Carlo standard error
in parentheses; multiply by two for executions. Full rate intervals and
all eight scenarios are in the supplied CSV.}
\label{tab:ustat_power}
\begin{tabular}{lrrrr}\toprule
 & \multicolumn{2}{c}{Efficiency gain} & \multicolumn{2}{c}{Joint gain}\\
Rule & Deploy (\%) & Use (MCSE) & Deploy (\%) & Use (MCSE)\\\midrule
All-pairs fixed Wald & 99.95 & 10000 (0) & 100.00 & 10000 (0) \\
All-pairs OBF & 99.95 & 4469 (29) & 100.00 & 2398 (12) \\
All-pairs Pocock & 99.75 & 3114 (38) & 100.00 & 1233 (10) \\
All-pairs HSD & 99.95 & 3932 (37) & 100.00 & 1498 (14) \\
Projection Gaussian & 98.15 & 3359 (51) & 100.00 & 715 (11) \\
Disjoint OBF & 99.95 & 4472 (29) & 100.00 & 2393 (12) \\
Disjoint betting & 96.60 & 4362 (52) & 100.00 & 1427 (15) \\
\bottomrule\end{tabular}\end{table}

All-pairs estimation reduces asymptotic variance in these settings,
but differences in boundaries, looks and calibration also affect stopping
and power. Identical raw component point estimates do not imply identical
estimated variances or stopping times. Table~\ref{tab:ustat_power} is a
descriptive comparison, not a paired significance or equivalence test.
All four calibration scenarios, all seven component procedures and both
crossing conventions are retained in the numerical supplement.
